\chapter{Connectivité des plateformes de titres à revenu fixe et de RFQ}\label{ch:nw:fixed-income-and-rfq-platform-connectivity}

Sur une plateforme de demande de cotation (RFQ), le \omterm{def:nw:fixed-income-and-rfq-platform-connectivity:timer}{temporisateur de cotation} du client court pendant des secondes, non des microsecondes~; pourtant les dealers répondent en
millisecondes. Une réponse qui arrive après le temporisateur n'est pas un moins bon prix mais aucun prix du tout, et certains clients n'attendent plus
attendre le temporisateur~: leurs règles exécutent la meilleure des toutes premières réponses, ou la première acceptable. Selon ces règles, un
dealer, plus lent de vingt millisecondes qu'un concurrent, perd des transactions qu'il aurait remportées sur le prix, et le modèle du chapitre établit que
diviser par deux le temps de tarification d'un auto-responder vaut près de quinze fois plus que de proposer un centime de mieux.

Le livre~2 a présenté la négociation par demande de cotation, les plateformes dealer-to-client, les courtiers interdealers et la négociation all-to-all~;
le livre~11 a construit l'auto-quoter avec son modèle de probabilité de gain et son prix de couverture (cover price)~; le livre~13 a traité la session FIX et le budget de latence.
Ce chapitre les relie~: les interfaces des plateformes, le pipeline de l'auto-responder, de la demande à la réponse, la latence qui
compte lorsque l'horloge est un \omterm{def:nw:fixed-income-and-rfq-platform-connectivity:timer}{temporisateur de cotation}, et les carnets d'ordres interdealers, qui sont de fait des bourses sans en porter le nom et se trouvent dans les
centres de données des chapitres~10 et~11. L'enchère est une simulation étiquetée, reposant sur le modèle de prix du livre~11.

\section{Interfaces des plateformes}

Un dealer se connecte à chaque plateforme dealer-to-client par une session FIX ou par l'interface propre à la plateforme, reçoit les demandes,
et y répond par des prix fermes valables pendant une durée indiquée. Les protocoles des plateformes diffèrent par ce que voit le client et par le temps pendant lequel
le prix du dealer engage celui-ci (\cref{dat:nw:fixed-income-and-rfq-platform-connectivity:platforms})~: une demande unilatérale, une demande bilatérale
de type request for market qui masque le côté du client, une request for stream dans laquelle les dealers actualisent leurs prix toutes les 250
millisecondes, ou un prix assorti d'un «~wire time~» pendant lequel il peut être frappé.

\begin{definition}[Temporisateur de cotation]\label{def:nw:fixed-income-and-rfq-platform-connectivity:timer}
Un \emph{temporisateur de cotation}\index{temporisateur de cotation} (quote timer) est la période, fixée par la plateforme ou par le client pour chaque demande de cotation, pendant
laquelle les réponses des dealers sont recueillies et restent exécutables~; une réponse qui arrive après son expiration n'est pas affichée, et un prix
affiché pendant cette période engage le dealer jusqu'à ce que le client exécute, passe ou que le temporisateur expire.
\end{definition}

\begin{dated}{2026-09}{Protocoles de demande de cotation et automatisation, d'après les plateformes}\label{dat:nw:fixed-income-and-rfq-platform-connectivity:platforms}
\begin{itemize}
\item Protocoles de TW SEF (avis de 2015)~: une RFQ est «~une enchère en temps réel avec plusieurs destinataires de RFQ~» dans laquelle le
  demandeur choisit le meilleur prix~; au moins trois destinataires non affiliés pour les transactions obligatoires~; request for market~; RFQ+
  à côté du carnet d'ordres~; request for stream, avec des prix actualisés «~toutes les 250 millisecondes~»~; sinon un prix assorti d'un «~wire
  time~».
\item L'AiEX de Tradeweb exécute les ordres issus du système de gestion des ordres (OMS) du client selon des règles prédéfinies~; un client peut «~exécuter au
  meilleur prix dès qu'il a reçu cinq réponses~» ou attendre un dealer particulier (2018).
\item MarketAxess (juin 2023)~: les protocoles d'automatisation représentaient «~près de 10~\% du volume total de négociation~» depuis le début de l'année~; son Adaptive
  Auto-X était en phase pilote.
\end{itemize}
\end{dated}

\section{Fonctionnement de l'auto-responder~: le pipeline de la demande à la réponse}

Un auto-responder est un pipeline~: la demande arrive par le réseau et le moteur FIX~; le pricer valorise l'obligation (la fair value du livre~11 à partir des
cours périmés, de la courbe de l'émetteur et des comparables)~; les limites et le risque sont contrôlés~; la réponse est renvoyée. Chaque étape
a une distribution de latence, et les demandes que l'auto-responder n'est pas autorisé à traiter, pour cause de taille, de liquidité ou de risque, vont à
une personne, qui répond en quelques secondes.

\omcode{code/firm/rfqlink/firm_rfqlink.py}{47}{53}{Temps d'une réponse~: les temps lognormaux des étapes, plus celui d'une personne pour la part des demandes hors des limites de l'auto-responder.}\label{lst:nw:fixed-income-and-rfq-platform-connectivity:response}

Le dealer du chapitre compte cinq étapes (réseau à l'entrée et à la sortie, \qty{0.5}{\milli\second} chacun, moteur FIX \qty{0.1}{\milli\second},
tarification \qty{20}{\milli\second}, limites \qty{1}{\milli\second}, médianes assorties de dispersions lognormales) et envoie 2~\% des demandes à
une personne, dont les réponses prennent six secondes en médiane. Ses quatre concurrents sont identiques, sauf une étape de tarification de
\qty{50}{\milli\second}. Tout cela relève d'hypothèses, dans les ordres de grandeur réalistes pour un auto-responder (livre~13).

\begin{definition}[Règle d'exécution automatique]\label{def:nw:fixed-income-and-rfq-platform-connectivity:rule}
Une \emph{règle d'exécution automatique}\index{règle d'exécution automatique} (automated execution rule) est une instruction préréglée du client, exécutée par la plateforme, qui
exécute une demande de cotation sans décision humaine~: à la meilleure réponse à l'expiration du temporisateur, à la meilleure des premières
$k$ réponses, à la première réponse située dans une tolérance autour d'un prix de référence, ou auprès d'un dealer désigné.
\end{definition}

\section{La latence qui compte lorsque l'horloge est un temporisateur de cotation}

Deux latences comptent. La première est de savoir si la réponse arrive avant le temporisateur~: avec un temporisateur de deux secondes, les étapes automatiques
ne manquent jamais, et la probabilité de manquer est presque exactement le produit de la part envoyée à une personne par la probabilité qu'une personne mette plus de deux secondes,
soit 1.94~\% pour le dealer du chapitre (voir la \cref{fig:nw:fixed-income-and-rfq-platform-connectivity:miss}). La seconde est de savoir si la réponse
arrive avant les réponses que la règle du client retiendra, ce qui dépend de cette règle.

\begin{omfigure}
\begin{tikzpicture}
\begin{axis}[width=11cm,height=5.6cm,xmode=log,ymode=log,xmin=100,xmax=31623,ymin=1e-5,ymax=0.1,
  xlabel={temporisateur de cotation (ms, log)},ylabel={probabilité de le manquer (log)},tick label style={font=\footnotesize},
  label style={font=\small},grid=major,grid style={gray!25},legend style={font=\footnotesize,at={(0.5,-0.3)},anchor=north,legend columns=2},
  unbounded coords=discard,table/col sep=comma]
\addplot[omThm,thick,mark=*,mark size=1.3pt] table[x=timer_ms,y=ours]{figdata/networks/25-fixed-income-and-rfq-platform-connectivity/miss.csv};
\addplot[omDef,thick,mark=square*,mark size=1.3pt] table[x=timer_ms,y=auto]{figdata/networks/25-fixed-income-and-rfq-platform-connectivity/miss.csv};
\legend{2~\% traitées par une personne, entièrement automatique}
\end{axis}
\end{tikzpicture}

\omcaption{Simulation~: probabilité que le dealer du chapitre réponde après le \omterm{def:nw:fixed-income-and-rfq-platform-connectivity:timer}{temporisateur de cotation}, en fonction de la durée de celui-ci. Entièrement
automatique, il ne manque jamais un temporisateur d'une seconde ou plus~; avec 2~\% des demandes traitées par une personne (six secondes en médiane), il
manque 1.94~\% des temporisateurs de deux secondes, et le palier vient de la personne, non du réseau. Données~: \texttt{fig\_rfq.py},
\texttt{nw\_rfq.miss\_curves()}.}\label{fig:nw:fixed-income-and-rfq-platform-connectivity:miss}
\end{omfigure}

La règle du client décide du reste. L'enchère tire le cours acheteur de chaque dealer comme le fait le livre~11~: l'estimation bruitée de la valeur de l'obligation par le dealer
diminuée d'une marge de 15 cents pour 100 de nominal, et le client ne vend qu'à un prix supérieur ou égal à son prix de réservation~; avec toutes les réponses reçues à temps
et la règle «~meilleure à l'expiration~», la simulation est exactement l'enchère du livre~11 (les tests le vérifient). Alors les règles divergent.

\omcode{code/firm/rfqlink/firm_rfqlink.py}{68}{85}{Les trois règles du client, et les demandes perdues faute de rapidité~: notre réponse était la meilleure acceptable dans les délais, et la règle en a choisi une autre.}\label{lst:nw:fixed-income-and-rfq-platform-connectivity:rules}

\begin{proposition}[Quand la vitesse ne vaut rien]\label{prop:nw:fixed-income-and-rfq-platform-connectivity:speed}
Avec la règle «~meilleure à l'expiration~», le taux de gain d'un dealer ne dépend de son temps de réponse que par la probabilité de manquer le temporisateur~;
avec «~meilleure des $k$ premières~» et «~première acceptable~», il baisse avec le rang du dealer dans l'ordre d'arrivée, si bien qu'une réponse plus rapide au
même prix gagne davantage.
\end{proposition}

\begin{proof}
À l'expiration, le client compare toutes les réponses arrivées à temps, quel que soit leur ordre. Avec les autres règles, une réponse hors des
$k$ premières, ou derrière une réponse acceptable antérieure, n'est jamais comparée, et un pipeline plus rapide accroît sa chance d'en faire partie.
\end{proof}

\begin{omfigure}
\begin{tikzpicture}
\begin{axis}[width=11cm,height=6cm,xmode=log,xmin=1,xmax=1000,ymin=0,ymax=0.65,
  xlabel={médiane de notre étape de tarification (ms, log)~; concurrents à 50 ms},ylabel={part des demandes que nous gagnons},
  tick label style={font=\footnotesize},label style={font=\small},grid=major,grid style={gray!25},
  legend style={font=\footnotesize,at={(0.5,-0.3)},anchor=north,legend columns=3},table/col sep=comma]
\addplot[omDef,thick,mark=*,mark size=1.2pt] table[x=pricing_ms,y=expiry]{figdata/networks/25-fixed-income-and-rfq-platform-connectivity/wins.csv};
\addplot[omMeth,thick,mark=square*,mark size=1.2pt] table[x=pricing_ms,y=first_k]{figdata/networks/25-fixed-income-and-rfq-platform-connectivity/wins.csv};
\addplot[omThm,thick,mark=triangle*,mark size=1.6pt] table[x=pricing_ms,y=first_ok]{figdata/networks/25-fixed-income-and-rfq-platform-connectivity/wins.csv};
\legend{meilleure à l'expiration,meilleure des trois premières,première acceptable}
\end{axis}
\end{tikzpicture}

\omcaption{Simulation~: la part des demandes que le dealer remporte parmi cinq, en fonction de son temps de tarification, selon trois règles des clients, tous les
dealers à la même marge. À l'échéance, la courbe reste plate tant que la queue de la distribution du temps de tarification atteint le temporisateur, à plusieurs centaines de millisecondes~; avec les deux autres règles, chaque
milliseconde compte. Données~: \texttt{fig\_rfq.py}, \texttt{nw\_rfq.curves()}.}\label{fig:nw:fixed-income-and-rfq-platform-connectivity:wins}
\end{omfigure}

À son étape de tarification de \qty{20}{\milli\second}, le dealer remporte 19.2~\% des demandes lorsque les clients retiennent la meilleure réponse à l'échéance, 26.4~\%
lorsqu'ils retiennent la meilleure des trois premières, et 39.0~\% lorsqu'ils retiennent la première réponse acceptable~: être plus rapide que ses quatre rivaux vaut
deux fois sa part. L'autre face de ces mêmes règles est le manque à gagner lié à la vitesse~: sous «~première acceptable~», 5.7~\% de toutes les demandes
sont celles où la réponse du dealer était la meilleure réponse acceptable à temps, mais où celle d'un rival plus rapide a été retenue~; sous «~meilleure des trois premières~»,
1.6~\%.

\begin{table}[ht]
\centering\small
\begin{tabular}{@{}>{\raggedright\arraybackslash}p{4.0cm}>{\centering\arraybackslash}p{2.6cm}>{\centering\arraybackslash}p{4.2cm}>{\centering\arraybackslash}p{2.6cm}@{}}
\toprule
Règle des clients & Taux de gain & Division par deux du temps de tarification (10 ms) & Un centime de mieux \\
\midrule
Meilleure à l'échéance & 19.2~\% & $+0.0$ points & $+0.9$ points \\
Meilleure des trois premières & 26.4~\% & $+2.3$ points & $+1.0$ points \\
Première acceptable & 39.0~\% & $+12.6$ points & $+0.9$ points \\
\bottomrule
\end{tabular}
\caption{Simulation~: la hausse du taux de gain obtenue en divisant par deux le temps de tarification, par rapport à une cotation meilleure d'un centime pour 100 de nominal (un point de
base de prix), selon chaque règle des clients. Données~: \texttt{nw\_rfq.values()}.}\label{tab:nw:fixed-income-and-rfq-platform-connectivity:value}
\end{table}

La valeur d'une milliseconde et celle d'un point de base ne sont donc pas comparables en général, mais seulement règle par règle
(le \cref{tab:nw:fixed-income-and-rfq-platform-connectivity:value}). Un dealer dont les clients attendent le temporisateur doit consacrer ses efforts au prix~; celui
dont les clients exécutent la première réponse acceptable doit d'abord consacrer ses efforts à son pipeline, puisque dix millisecondes valent presque quinze centimes de
prix. Les règles retenues par ses clients sont une information que le dealer peut estimer à partir de ses propres gains et pertes selon le temps de réponse.

\section{Les carnets d'ordres interdealers}

Les marchés interdealers d'obligations d'État sont des carnets d'ordres à cours limité centralisés, gérés comme des bourses et situés sur les mêmes campus que
les marchés de contrats à terme et d'options des chapitres~10 et~11 (l'\cref{dat:nw:fixed-income-and-rfq-platform-connectivity:idb}). Les marchés américains de BrokerTec
s'apparient chez CME, au NY5 de Secaucus, et ses marchés européens au LD4.2 de Slough, chacun doté d'un site de reprise~; BrokerTec Chicago, un
second carnet d'ordres pour les Treasuries les plus activement négociées, est installé au centre de données de CME à Aurora, à côté des futures sur Treasuries, pour négocier le cash
contre les futures. Sur ces carnets, tout ce qui est vu aux chapitres~1 à~9 s'applique~: \omterm{def:nw:colocation-products-and-how-they-are-sold:colo}{colocation}, \omterm{def:nw:colocation-products-and-how-they-are-sold:mmr}{interconnexions} (cross-connects), feeds, passerelles~; les
secondes des plateformes RFQ redeviennent des microsecondes.

\begin{dated}{2026-09}{Carnets d'ordres interdealers d'obligations d'État sur CME Globex, d'après la documentation de CME}\label{dat:nw:fixed-income-and-rfq-platform-connectivity:idb}
Marchés américains de BrokerTec~: production à NY5, reprise à Aurora~; marchés européens~: production à LD4.2, reprise à NY5~; objectifs de délai de reprise
de quatre heures (États-Unis) et de deux heures (Europe). BrokerTec Chicago (carnet d'ordres secondaire des Treasuries américaines les plus actives)~: le centre de données CME Globex à
Aurora, Illinois.
\end{dated}

La distance entre NY5 et Aurora est celle du basis trade cash--futures~: les futures sur Treasuries américains s'apparient à Aurora, le
carnet cash principal à Secaucus, et les deux sont séparés par le couloir New Jersey--Chicago du chapitre~10, la route sur laquelle la fibre et les faisceaux hertziens
se disputent la course. Un second carnet cash à Aurora réunit les deux jambes dans un même bâtiment pour les firmes qui les négocient ensemble.

\begin{method}[Connecter un desk obligataire]\label{met:nw:fixed-income-and-rfq-platform-connectivity:desk}
\begin{enumerate}
\item Dresser la liste des plateformes et des carnets que le desk négocie~; pour chacun, le protocole, le type de session, le temporisateur ou le temps de transmission, et le centre de
  données.
\item Construire l'auto-répondeur sous forme de pipeline avec un budget de latence par étape~; mesurer la part envoyée à une personne et les temps de réponse
  de cette personne.
\item Estimer les règles d'exécution des clients à partir des gains et des pertes selon le temps de réponse~; consacrer ses efforts à la vitesse là où les règles la récompensent, et au prix
  là où elles ne la récompensent pas.
\item Recourir à la \omterm{def:nw:colocation-products-and-how-they-are-sold:colo}{colocation} pour les carnets d'ordres, comme pour toute bourse~; pour le cash contre les futures, choisir entre deux sites ou un seul.
\item Surveiller les dépassements du temporisateur, les réponses postérieures aux $k$ premières, et les gains perdus faute de vitesse, par plateforme et par client.
\end{enumerate}
\end{method}

\section{Tutoriel~: deux secondes pour répondre}\label{tut:nw:fixed-income-and-rfq-platform-connectivity:tut}

\begin{tutorial}
\textbf{Objectif.} Simuler le pipeline d'un auto-répondeur, le temporisateur et trois règles des clients sur l'adjudication du livre~11.
\textbf{État final~:} la \cref{fig:nw:fixed-income-and-rfq-platform-connectivity:wins} et
le \cref{tab:nw:fixed-income-and-rfq-platform-connectivity:value}.
\begin{tutsteps}
\item \textbf{Pipelines.} \texttt{firm\_rfqlink.Pipeline} et \texttt{Stage}~; \texttt{response\_ms}
  (le \cref{lst:nw:fixed-income-and-rfq-platform-connectivity:response}) et \texttt{miss\_prob} pour le temporisateur.
\item \textbf{Adjudication.} \texttt{auction} avec chaque règle (le \cref{lst:nw:fixed-income-and-rfq-platform-connectivity:rules}), contrôlé
  par rapport à \texttt{firm\_rfqmm.simulate} lorsque toutes les réponses arrivent à temps.
\item \textbf{Valeur.} \texttt{value\_of} pour un pipeline plus rapide et un meilleur prix~; \texttt{nw\_rfq.curves} sur l'ensemble des temps de tarification.
\end{tutsteps}
\textbf{Pour aller plus loin.} Laisser chaque client choisir sa règle au hasard, avec des parts estimées à partir des données, et laisser la marge réagir~:
l'équilibre du livre~11 sous la règle de la première réponse acceptable.
\end{tutorial}

\section{Construction~: le modèle de connectivité RFQ}\label{bld:nw:fixed-income-and-rfq-platform-connectivity:rfqlink}

\begin{build}
\bfield{Rôle} Temporisateurs, pipelines et règles des clients pour un dealer sur les plateformes RFQ~: les lignes taux du plan du chapitre~29.
\bfield{Interface} \texttt{firm\_rfqlink}~: \texttt{Stage}, \texttt{Pipeline}, \texttt{response\_ms}, \texttt{miss\_prob},
\texttt{auction}, \texttt{value\_of}, \texttt{scaled}, \texttt{log\_grid}.
\bfield{Règles} Les prix suivent le modèle \texttt{firm.rfqmm} du livre~11, inchangé~; les paramètres du pipeline et des clients sont des hypothèses énoncées~;
les résultats sont une simulation étiquetée.
\bfield{Tests d'acceptation} \texttt{code/firm/rfqlink/tests/}~: identité avec \texttt{firm\_rfqmm.simulate} lorsque toutes les réponses arrivent à temps,
pipelines et temporisateurs déterministes, et l'ordre des règles.
\bfield{Pour aller plus loin} Temps de transmission et last look sur les flux~; répartition des règles des clients estimée à partir des données~; limiteurs de débit sur les sessions.
\end{build}

\begin{omsources}
\item TW SEF LLC, Trading and Execution Protocols (2015).
\item The TRADE, «~Automating trade execution, intelligently~» (2018); Markets Media on MarketAxess's Adaptive Auto-X (2023).
\item CME Group Client Systems Wiki~: BrokerTec disaster recovery process and scenarios~; BrokerTec Chicago.
\end{omsources}

\section{Exercices}

\begin{exercise}[$\star$]\label{exo:nw:fixed-income-and-rfq-platform-connectivity:1}
Expliquer le sort d'une réponse parvenue après le \omterm{def:nw:fixed-income-and-rfq-platform-connectivity:timer}{temporisateur de cotation}.
\end{exercise}

\begin{exercise}[$\star$]\label{exo:nw:fixed-income-and-rfq-platform-connectivity:2}
Indiquer où les carnets d'ordres américains et européens de BrokerTec s'apparient, et où se situent leurs sites de reprise.
\end{exercise}

\begin{exercise}[$\star$]\label{exo:nw:fixed-income-and-rfq-platform-connectivity:3}
Expliquer pourquoi la probabilité de manquer un temporisateur de deux secondes est presque égale à la part des demandes envoyées à une personne.
\end{exercise}

\begin{exercise}[$\star\star$]\label{exo:nw:fixed-income-and-rfq-platform-connectivity:4}
À partir de la \cref{prop:nw:fixed-income-and-rfq-platform-connectivity:speed}, pourquoi diviser par deux le temps de pricing ne rapporte-t-il rien sous «~meilleure à
l'échéance~»~?
\end{exercise}

\begin{exercise}[$\star\star$]\label{exo:nw:fixed-income-and-rfq-platform-connectivity:5}
Un client exécute «~au meilleur prix dès qu'il dispose de cinq réponses~» et interroge cinq dealers. Laquelle des règles du chapitre correspond-elle à ce cas~?
\end{exercise}

\begin{exercise}[$\star\star$]\label{exo:nw:fixed-income-and-rfq-platform-connectivity:6}
Pourquoi un trader de basis cash--futures pourrait-il préférer un deuxième carnet d'ordres Treasury chez Aurora~?
\end{exercise}

\begin{exercise}[$\star\star\star$]\label{exo:nw:fixed-income-and-rfq-platform-connectivity:7}
\emph{Programmation.} Avec \texttt{nw\_rfq.curves}, à quel temps de pricing le taux de gain du dealer sous «~première acceptable~» tombe-t-il
sous son taux sous «~meilleure à l'échéance~»~?
\end{exercise}

\begin{exercise}[$\star\star\star$]\label{exo:nw:fixed-income-and-rfq-platform-connectivity:8}
\emph{Trouver la faille.} «~Nos clients disposent de temporisateurs de deux secondes, donc notre pricer de 20~millisecondes est cent fois plus rapide qu'il ne doit l'être.~»
\end{exercise}

\section{Problème~: Deux secondes pour répondre}

\begin{problem}[{Problème du week-end --- temporisateurs, règles et pipeline d'un dealer}]\label{pb:nw:fixed-income-and-rfq-platform-connectivity:1}
Un dealer répond aux demandes de cotation (RFQ) sur des obligations d'entreprise face à quatre concurrents. Son pipeline et celui de ses concurrents sont ceux du chapitre~; le
\omterm{def:nw:fixed-income-and-rfq-platform-connectivity:timer}{temporisateur de cotation} est de deux secondes.

\medskip\noindent\textbf{Partie I --- Le temporisateur.}
\begin{enumerate}
\item Donner la probabilité que le dealer manque le temporisateur, et expliquer d'où elle provient.
\item Donner la probabilité de manquer le temporisateur si le pipeline était entièrement automatique.
\item Déterminer à partir de quelle durée de temporisateur le pipeline entièrement automatique ne manque plus aucune réponse.
\item Comparer le coût d'une réponse manquée pour le dealer à celui d'une réponse perdante.
\end{enumerate}

\medskip\noindent\textbf{Partie II --- Les règles.}
\begin{enumerate}[resume]
\item Déterminer la part des demandes que le dealer remporte sous chaque règle.
\item Déterminer la part perdue à cause de la vitesse sous chaque règle.
\item Expliquer pourquoi elle est nulle à l'échéance.
\item Expliquer pourquoi le dealer remporte davantage sous «~première acceptable~» qu'à l'échéance.
\end{enumerate}

\medskip\noindent\textbf{Partie III --- La vitesse contre le prix.}
\begin{enumerate}[resume]
\item Déterminer ce que rapporte la division par deux du temps de pricing sous chaque règle.
\item Déterminer le gain procuré par une cotation meilleure d'un centime.
\item Calculer combien de centimes de prix valent dix millisecondes sous «~première acceptable~».
\item Dire comment le dealer peut savoir quelles règles utilisent ses clients.
\end{enumerate}

\medskip\noindent\textbf{Partie IV --- Le verdict.}
\begin{enumerate}[resume]
\item Énoncer le \emph{résultat nommé}~: la probabilité de manquer le \omterm{def:nw:fixed-income-and-rfq-platform-connectivity:timer}{temporisateur de cotation} pour la distribution de latence du pipeline, et la part
  d'activité perdue au profit d'un dealer plus rapide sous une règle d'exécution «~première acceptable~».
\item Dire quelles données sont publiées et lesquelles sont supposées.
\item Dire où le dealer doit consacrer son prochain mois d'ingénierie.
\item Analyser ce qui se passe si chaque concurrent divise lui aussi son temps de pricing par deux.
\item Dire comment gérer la part envoyée à une personne.
\item Dire ce qui change dans les carnets d'ordres inter-dealers.
\item Dire où cela s'insère dans le plan du chapitre~29.
\item Résumer en une phrase le moment où une milliseconde compte sur une plateforme dont le temporisateur est de deux secondes.
\end{enumerate}
\end{problem}

\section{Questions d'entretien}

\begin{interviewq}[$\star$ \iqroles{développeur}]\label{iq:nw:fixed-income-and-rfq-platform-connectivity:1}
Quelles sont les étapes d'un auto-responder RFQ, et laquelle vous semble la plus lente~?
\end{interviewq}

\begin{interviewq}[$\star\star$ \iqroles{trader}]\label{iq:nw:fixed-income-and-rfq-platform-connectivity:2}
Pourquoi un client exécuterait-il la première réponse acceptable plutôt que d'attendre la meilleure~?
\end{interviewq}

\begin{interviewq}[$\star\star$ \iqroles{développeur, chercheur}]\label{iq:nw:fixed-income-and-rfq-platform-connectivity:3}
Comment estimeriez-vous si votre temps de réponse vous fait perdre des affaires sur une plateforme~?
\end{interviewq}

\begin{interviewq}[$\star\star$ \iqroles{développeur}]\label{iq:nw:fixed-income-and-rfq-platform-connectivity:4}
La latence p99 de votre auto-responder a doublé après une mise en production. Que vérifiez-vous~?
\end{interviewq}

\begin{interviewq}[$\star\star$ \iqroles{chercheur}]\label{iq:nw:fixed-income-and-rfq-platform-connectivity:5}
Un dealer doit-il coter plus serré ou répondre plus vite~? Comment trancheriez-vous~?
\end{interviewq}

\begin{interviewq}[$\star\star\star$ \iqroles{développeur, chercheur}]\label{iq:nw:fixed-income-and-rfq-platform-connectivity:6}
Concevez la connectivité et l'auto-responder d'un dealer qui répond sur trois plateformes obligataires et négocie sur un carnet inter-dealers.
\end{interviewq}
